% !TEX root = ../Thesis.tex
\section{Five Estimators on the Cohort}

In this chapter we apply the machinery the risk estimation to the MASTER DAPT cohort
and reports the findings.
The main point is answering four questions:

\begin{itemize}
  \item Is there evidence of risk difference in the Treatments effects across patients?
  \item If yes, does it apply best to some patients and not to others?
  \item Can we derive decisions from the estimated heterogeneity that improve on the best fixed arm?
  \item How large of a Variation would have been needed to be detected by the protocol?
\end{itemize}

\paragraph{The five estimators.}
Five effect estimators are fitted on the same 4,579 patients, the same
covariate set and the same five endpoints of Chapter 4, each belonging to a
different methodological family:

\begin{itemize}
  \item a \textbf{two-model meta-learner} (T-learner), the difference of the two
        calibrated per-arm random forests of Section~\ref{sec:per-arm};
  \item a \textbf{causal forest} (\texttt{CausalForestDML}), which targets the
        effect directly through double machine learning with honest splitting,
        with gradient-boosted first stages for both nuisance functions
        $\mathbb{E}[Y \mid x]$ and $\mathbb{E}[T \mid x]$;
  \item a \textbf{Bayesian causal forest} (BCF), which places separate priors and
        separate tree ensembles on the prognostic surface $\mu(x)$ and on the
        effect surface $\tau(x)$, fixes the propensity at the known $0.5$, and
        returns a full posterior rather than a point estimate;
  \item an \textbf{in-context model} (CausalPFN), a transformer pre-trained on
        synthetic causal problems that conditions on the cohort at inference
        time and is never fitted to it in the usual sense;
  \item an \textbf{interaction forest} (S-learner), a single classification
        forest on the augmented design $[\,X,\;T,\;X \!\cdot\! T\,]$, whose
        implied CATE is obtained by toggling the treatment column.
\end{itemize}

None of the five is included because it is the best estimator in the abstract.
They are included because their failure modes are different: the T-learner
inherits the errors of two outcome models, the causal forest can only find
structure its splitting criterion can express, the BCF can be dominated by its
own regularization, the in-context model carries a domain bias from its
pre-training distribution rather than sampling variance from this cohort, and
the S-learner pools the two arms and shrinks weak interactions toward zero. A
result common to all five is therefore hard to attribute to any one of them.

\paragraph{The common protocol.}
All estimators use the same preprocessing pipeline, with k-NN imputation and
robust scaling fitted only on the training data. Heterogeneity is evaluated
by fitting the model on 70\% of the data, ranking the held-out 30\% by
predicted CATE, and applying EconML's doubly-robust tester with the known
randomization propensity. Performance is assessed using the BLP slope,
AUTOC, and QINI. AUTOC emphasizes the top of the ranking, whereas QINI
weights groups by size, so both are reported to ensure that negative
heterogeneity results are robust to the choice of weighting.

\section{The Positive Control: The Average Effect Is Recovered}
\label{sec:positive-control}

Before asking whether the effect varies, one has to establish that the
estimators see the effect at all. Table~\ref{tab:ate-recovery} compares the
trial's unadjusted arm difference with the mean CATE produced by each of the
five estimators, endpoint by endpoint. Under the convention of Chapter 4 a
positive value is the benefit of abbreviating.

\begin{table}[h]
\centering
\caption{Average effect recovery. Raw ATE is the unadjusted difference in event
rate between arms; the remaining columns are the cohort mean of each
estimator's per-patient CATE. Positive $=$ abbreviating lowers the event.}
\label{tab:ate-recovery}
\small
\setlength{\tabcolsep}{4pt}
\begin{tabular}{lrrrrrr}
\toprule
Endpoint & Raw ATE & T-learner & Causal for. & BCF & CausalPFN & Interact.\ for. \\
\midrule
Bleeding             & $+0.045$ & $+0.055$ & $+0.047$ & $+0.047$ & $+0.043$ & $+0.071$ \\
Bleeding BARC 2/3/5  & $+0.028$ & $+0.034$ & $+0.029$ & $+0.029$ & $+0.021$ & $+0.052$ \\
MI                   & $-0.005$ & $-0.004$ & $-0.004$ & $-0.003$ & $-0.003$ & $-0.029$ \\
Cardiovascular death & $+0.003$ & $+0.002$ & $+0.003$ & $+0.003$ & $-0.005$ & $+0.001$ \\
Stroke               & $+0.005$ & $+0.005$ & $+0.004$ & $+0.005$ & $+0.001$ & $+0.020$ \\
\bottomrule
\end{tabular}
\end{table}

The agreement is close and, on the two bleeding endpoints, close in the way
that matters: prolonging therapy costs roughly four to five percentage points
of bleeding risk, which is the trial's own published conclusion, reached here
by five estimators that share no code path beyond the preprocessing. 

On the ischemic side, all estimators return effects within 0.5 percentage
points of zero, with MI consistently favouring prolonged therapy. The BCF
also gives credible intervals: $[+0.030,\,+0.067]$ for bleeding, excluding
zero, compared with $[-0.011,\,+0.005]$ for MI and $[-0.004,\,+0.010]$ for
cardiovascular death, both spanning zero.

The interaction forest is the main outlier, overstating all effects and MI
by a factor of six. This is expected because its in-sample CATE estimates
come from an outcome-prediction model and can absorb prognostic differences.
Consequently, the S-learner's dispersion can be viewed as an \emph{upper
bound} on genuine heterogeneity: if even this bound is uninformative, the
evidence for meaningful heterogeneity is even weaker.

\section{The Negative Result: The Dispersion Is Noise}

Every estimator produces a distribution of per-patient effects, and every such
distribution has non-zero width. That width, on its own, means nothing: a
constant true effect estimated with error also produces a spread of estimates,
and the width of that spread grows with the flexibility of the model, not with
the amount of heterogeneity in the data. The question is therefore never
whether the dispersion is positive, but whether it exceeds what noise alone
would produce, and whether the resulting \emph{ranking} of patients tracks
anything real.

\paragraph{Dispersion, and a noise floor to compare it against.}
Table~\ref{tab:cate-dispersion} reports the standard deviation of the
per-patient CATE for each estimator. The five columns disagree by an order of
magnitude — from $0.003$ to $0.093$ on the same patients and the same endpoint
— which is itself the first indication that what is being measured is a
property of the estimators rather than of the cohort. If the spread reflected
real variation in treatment response, it would not depend so strongly on who
is measuring it.

\begin{table}[h]
\centering
\caption{Standard deviation of the per-patient CATE by estimator. The BCF
column reports the posterior mean of $\mathrm{sd}(\tau(x))$, and the placebo
column the same quantity obtained after permuting the treatment labels, i.e.\
the dispersion produced by estimation noise when the true effect is exactly
constant. The floor is computed only for the two endpoints with enough events
to support it.}
\label{tab:cate-dispersion}
\small
\setlength{\tabcolsep}{4pt}
\begin{tabular}{lccccc|c}
\toprule
Endpoint & T-learner & Causal for. & BCF & CausalPFN & Interact.\ for. & Placebo floor \\
\midrule
Bleeding             & 0.037 & 0.006 & 0.039 & 0.093 & 0.090 & 0.035 \\
Bleeding BARC 2/3/5  & 0.032 & 0.004 & 0.027 & 0.059 & 0.085 & 0.023 \\
MI                   & 0.024 & 0.003 & 0.015 & 0.009 & 0.057 & --- \\
Cardiovascular death & 0.028 & 0.004 & 0.018 & 0.006 & 0.045 & --- \\
Stroke               & 0.020 & 0.007 & 0.014 & 0.006 & 0.046 & --- \\
\bottomrule
\end{tabular}
\end{table}

The last column is what turns the comparison into a test. Permuting the
treatment labels destroys the effect while leaving the covariates, the outcome
rates and the estimation procedure untouched, so refitting on permuted data
measures exactly how much apparent heterogeneity the method manufactures from
nothing. On bleeding the BCF's observed spread is $0.039$ against a floor of
$0.035$; on BARC 2/3/5 it is $0.027$ against $0.023$. The excess is $0.004$ in
both cases, and the posterior probability that the observed spread exceeds the
mean of the placebo distribution is $0.71$ and $0.76$ respectively — a coin
flip's distance from the noise, on the two endpoints where the cohort has the
most events to offer. Figure~\ref{fig:placebo-floor} shows the two
distributions overlapping almost completely.

\begin{figure}[h]
\centering
\includegraphics[width=1\textwidth]{CATE_placebo_floor.png}
\caption{Posterior of the cross-patient spread $\mathrm{sd}(\tau(x))$ under the
Bayesian causal forest (red) against the placebo distribution obtained by
permuting the treatment labels (grey), for the two best-powered endpoints.
Dashed lines mark the two means. A real heterogeneity signal would place the
red distribution clear of the grey one.}
\label{fig:placebo-floor}
\end{figure}

\paragraph{The ranking does not beat random targeting.}
The noise floor is a model-internal diagnostic: it compares a fitted spread
against a spread obtained under the null, without consulting the outcomes on a
held-out set. The decisive test is the other one. Table~\ref{tab:drtester}
reports the doubly-robust validation of the ranking on the held-out 30\% for
the three estimators that share the protocol exactly.

\begin{table}[h]
\centering
\caption{Doubly-robust validation of the CATE ranking on a held-out 30\% split.
AUTOC is the targeting gain over random with $1/q$ weighting; a positive value
with a small $p$ would indicate a usable ranking.}
\label{tab:drtester}
\begin{tabular}{lcccccc}
\toprule
& \multicolumn{2}{c}{Causal forest} & \multicolumn{2}{c}{T-learner} & \multicolumn{2}{c}{BCF} \\
\cmidrule(lr){2-3}\cmidrule(lr){4-5}\cmidrule(lr){6-7}
Endpoint & AUTOC & $p$ & AUTOC & $p$ & AUTOC & $p$ \\
\midrule
Bleeding             & $-0.013$ & 0.18 & $-0.015$ & 0.14 & $+0.021$ & 0.07 \\
Bleeding BARC 2/3/5  & $-0.005$ & 0.31 & $+0.003$ & 0.41 & $+0.015$ & 0.10 \\
MI                   & $-0.009$ & 0.10 & $-0.006$ & 0.24 & $-0.003$ & 0.34 \\
Cardiovascular death & $-0.000$ & 0.50 & $-0.006$ & 0.26 & $+0.006$ & 0.22 \\
Stroke               & $+0.005$ & 0.19 & $+0.004$ & 0.16 & $+0.006$ & 0.09 \\
\bottomrule
\end{tabular}
\end{table}

Not one of the fifteen entries reaches significance at the conventional 5\%
level — the closest is the BCF on bleeding, at $p = 0.07$ and with the sign
opposite to the causal forest's on the same endpoint — and the signs are not even
consistent across estimators on the same endpoint. The QINI weighting agrees:
under the causal forest the largest QINI is $-0.007$ on bleeding
($p = 0.08$) and every other value is within $0.002$ of zero, so the negative
result is not an artefact of a weighting that over-emphasises the head of the
ranking. The BLP slopes point the same way and are, if anything, blunter: on
every endpoint but stroke the regression of the doubly-robust effect score on
the predicted CATE has a \emph{negative} slope ($-3.16$ on bleeding, $p = 0.25$),
which is to say that following the model's ranking is, within noise, very
slightly worse than ignoring it. Calibration $R^2$ is negative on three of the
five endpoints and essentially zero on a fourth, meaning that the observed
group effects are predicted worse by the model's group predictions than by
their own overall mean.

The two remaining estimators are evaluated with weaker but independent
protocols and reach the same place. The in-context model's bootstrapped TOC
bands cover zero on every endpoint, and its strongest Spearman correlation
between any covariate and the predicted CATE is $0.21$ — no patient
characteristic meaningfully moves its estimate. The interaction forest points
the same way from inside its own design matrix: the treatment column carries
between $0.0009$ and $0.0018$ of the forest's total splitting importance
depending on the endpoint, so the ensemble that is free to split on treatment,
on covariates and on their product barely uses the treatment at all.

\begin{figure}[h]
\centering
\includegraphics[width=1\textwidth]{CATE_sorted_forest.png}
\caption{Per-patient CATE of the tuned causal forest, patients sorted by
predicted effect, one panel per endpoint. On both bleeding endpoints the whole
cohort sits on the same side of zero: the ordering exists, but the range it
spans is a fraction of the average effect.}
\label{fig:sorted-cate}
\end{figure}

\begin{figure}[h]
\centering
\includegraphics[width=1\textwidth]{CATE_rate_toc.png}
\caption{TOC curves for the tuned causal forest on the held-out split: gain
over random targeting as the treated fraction grows, with 95\% confidence
bands. A usable ranking would lift the curve clear of the dashed line at small
treated fractions, where the clinical decision is made. The bands cover zero
throughout on every endpoint.}
\label{fig:rate-toc}
\end{figure}

Figures~\ref{fig:sorted-cate} and \ref{fig:rate-toc} show the same result from
the two sides that matter. The sorted CATE curves are monotone by construction
and therefore always look like heterogeneity; what they actually show, on the
bleeding endpoints, is a cohort in which every single patient is estimated to
benefit from abbreviating, spanning $+0.032$ to $+0.068$ around a mean of
$+0.047$ — a total range smaller than the average effect itself. The TOC curves then ask whether that ordering buys anything,
and answer no: at every treated fraction, the doubly-robust gain over random
targeting is indistinguishable from zero.

\paragraph{What the two results say together.}
Read separately, either half could be dismissed. A positive control on its own
shows that the estimators can find an average effect, which is a much easier
task than finding heterogeneity. A null result on its own is compatible with
having searched badly. Together they are informative: the same estimators, the
same data, the same validation protocol recover an average effect that matches
the published trial and fail to find any variation around it above the level
that noise alone produces. The distinction this chapter is entitled to claim is
between \emph{nothing was found} and \emph{a search with tools that demonstrably
work found nothing} --- and Section~\ref{sec:delta} makes even that claim
quantitative by measuring how large a heterogeneity those tools would have
needed in order to find it.

\section{An Episode: The Model That Collapsed}

One episode from the estimation work is worth telling at length, not for its
result but for the reasoning it forced, which is the reasoning the whole chapter
depends on.

Following Chapter 2, the causal forest's hyperparameters are tuned on the same
criterion used to evaluate it: the cross-validated lower confidence bound of
the AUTOC on the bleeding endpoint, $\mathrm{AUTOC} - z \cdot \mathrm{SE}$
averaged over folds. The search is deliberately wide — depth unbounded, leaf
size from 1 to 200, the first-stage nuisance flexibility tuned jointly with the
forest — and it is run as a parallel Optuna study over a shared journal;
the configuration reported here was selected over 23,361 completed trials. Two
defences against the winner's curse are built into the objective: the lower
bound rather than the point estimate, so that a large-but-uncertain score
cannot win, and a hard zero for any configuration that returns a constant CATE,
so that a degenerate model cannot score well by having no ranking to be wrong
about.

That second defence exists because of what happened without it. A tuned
configuration returned $\mathrm{std}(\mathrm{CATE})$ of exactly zero — every
patient assigned the identical effect — and the immediate question was whether
this was a bug in the pipeline or an answer from the data. It was traced to a
single hyperparameter, the forest's \texttt{min\_impurity\_decrease}, the gate
that refuses a split whose improvement falls below a threshold: set high enough,
no covariate ever justifies a split of the effect surface, and the forest
returns the ATE for everyone. That is not a malfunction. It is the estimator
reporting, in the bluntest available form, that it cannot find a covariate
worth splitting on.

\paragraph{The flexibility ladder.}
Rather than settle the question with a single tuned model, three explicit
configurations of increasing capacity are fitted, scaling the forest and its
first-stage nuisance models together: \texttt{regularized} (depth 2, leaves of
at least 40), \texttt{medium} (depth 4, leaves of 15) and \texttt{flexible}
(depth 8, leaves of 5). Table~\ref{tab:flex-ladder} reports what happens.

\begin{table}[h]
\centering
\caption{Mean and standard deviation of the per-patient CATE across three
explicit flexibility levels and the tuned model. The mean is stable; only the
dispersion moves.}
\label{tab:flex-ladder}
\small
\begin{tabular}{lcccc}
\toprule
& \multicolumn{4}{c}{mean CATE / std CATE} \\
\cmidrule(lr){2-5}
Endpoint & Regularized & Medium & Flexible & Tuned \\
\midrule
Bleeding             & $0.046$ / $0.004$ & $0.045$ / $0.006$ & $0.046$ / $0.016$ & $0.047$ / $0.006$ \\
Bleeding BARC 2/3/5  & $0.028$ / $0.002$ & $0.026$ / $0.005$ & $0.026$ / $0.014$ & $0.029$ / $0.004$ \\
MI                   & $-0.004$ / $0.000$ & $-0.005$ / $0.004$ & $-0.005$ / $0.012$ & $-0.004$ / $0.003$ \\
Cardiovascular death & $0.003$ / $0.000$ & $0.003$ / $0.004$ & $0.004$ / $0.011$ & $0.003$ / $0.004$ \\
Stroke               & $0.004$ / $0.000$ & $0.004$ / $0.003$ & $0.004$ / $0.006$ & $0.004$ / $0.007$ \\
\bottomrule
\end{tabular}
\end{table}

The pattern is unambiguous and it is the pattern of a null. Across a factor of
four in tree depth and a factor of eight in minimum leaf size, the \emph{mean}
CATE does not move at all — bleeding stays at $0.045$--$0.047$ on every rung of
the ladder — while the \emph{dispersion} grows by a factor of four. Under the
regularized preset the three ischemic endpoints collapse to a constant: stroke
takes exactly one distinct value across all 4,579 patients, and cardiovascular
death and MI take three, spanning a range of $2 \times 10^{-5}$ and
$1 \times 10^{-4}$ respectively. Under the flexible preset the same endpoints
show dispersions of $0.011$ and $0.012$, comfortably larger than their own
average effects.

A quantity that is invariant to model capacity is a property of the data; a
quantity that scales with model capacity is a property of the model. The
average effect is the first, the heterogeneity is the second. This is why the
collapse was reported as an answer rather than patched away, and why the
flexible preset's apparent heterogeneity is not evidence against the
conclusion: those effect surfaces are fitted to noise, and
Table~\ref{tab:drtester} is where they are asked to prove otherwise on
held-out data, and fail.

\paragraph{A note on the tuning score itself.}
The best cross-validated AUTOC lower bound found by the search was $+0.014$ on
the bleeding endpoint — positive, and encouraging in isolation. The same model,
evaluated once on a held-out split by the doubly-robust tester, returns
$-0.013$ ($p = 0.18$). The maximum of a noisy criterion over tens of thousands
of configurations is itself an optimistic statistic, and the lower-bound
objective attenuates that optimism without eliminating it. It is a concrete
reason why no heterogeneity claim in this thesis rests on a tuning score, and
why the held-out doubly-robust evaluation is the only figure of merit the
conclusions are allowed to use.

\todo[inline]{The Optuna study has since been extended to 161,336 trials and
its current optimum differs from the configuration of the saved model reported
here. Either refit the tuned model from the current optimum or freeze the study
at the reported state before the final draft.}

\section{The Trade-off Plane}

Every result so far is per endpoint, and the clinical question is not. A
physician deciding whether to abbreviate is not trading a bleeding event
against nothing; they are trading a bleeding event against an ischemic one.
Because the per-endpoint forests give each patient a CATE for every outcome,
those trades can be looked at directly: put the bleeding benefit of
abbreviating on one axis and the ischemic benefit on the other, and each
patient becomes one point on a plane whose four quadrants have fixed clinical
meanings. Top-right is win--win, abbreviate; bottom-left is lose--lose, keep
prolonged; and \textbf{bottom-right} is the quadrant that would justify
personalization — abbreviating stops bleeds but costs ischemic events, so the
decision genuinely depends on the patient. The ischemic axis aggregates the
three ischemic CATEs with fixed weights ($0.4$ MI, $0.4$ stroke, $0.2$
cardiovascular death), the same weights used throughout the thesis.

\paragraph{A sign check before the plane is read.}
An axis is only usable if the model's mean CATE agrees with the trial's raw
arm difference on that endpoint; where it does not, the endpoint's effect
surface is unreliable and any structure read off it is an artefact of the fit.
The check is run per endpoint before anything else is plotted. Under the model
reported here all five endpoints pass — Table~\ref{tab:ate-recovery} is that
check in tabular form — and overall bleeding is retained as the bleeding axis
for its larger event count. The gate is not decorative: under an earlier
per-endpoint tuned forest the BARC 2/3/5 CATE mean carried the opposite sign to
its own ATE, and that model was discarded as an axis on exactly this ground.

\todo[inline]{The outline states that BARC 2/3/5 fails the sign check and must
be dropped as an axis. With the currently saved tuned forest it passes
(mean $+0.029$ against a raw ATE of $+0.028$). Confirm which model version the
final text should describe.}

\begin{figure}[h]
\centering
\includegraphics[width=1\textwidth]{CATE_tradeoff_plane.png}
\caption{The trade-off plane: bleeding benefit of abbreviating against the
weighted ischemic benefit, one point per patient. Left, patient density;
right, the same points tinted by the conflict score. The entire cohort sits in
the right half-plane, and the cloud straddles the horizontal axis without any
concentration in the bottom-right quadrant that a personalized rule would
require.}
\label{fig:tradeoff-plane}
\end{figure}

\paragraph{What the plane shows.}
Figure~\ref{fig:tradeoff-plane} is a single compact cloud. Its horizontal
extent runs from $+0.032$ to $+0.068$ — every patient on the abbreviate side,
none anywhere near zero — while vertically the cloud spans roughly $-0.004$ to
$+0.010$, straddling the ischemic axis in a way that reflects the near-null
ischemic effect rather than any grouping. There is no separated mass in the
bottom-right, and there is no visible gap anywhere that would let one draw a
boundary. Clustering the standardized plane confirms the impression rather than
overturning it: the best silhouette over $k = 2 \ldots 6$ selects three
clusters that are simply three slices of one blob, ordered along the ischemic
axis, with mean bleeding CATE of $0.043$, $0.047$ and $0.053$ — a spread
smaller than the bootstrap uncertainty of a single patient's position.

\paragraph{The plane with error bars.}
A single fit gives the plane no uncertainty, so it cannot distinguish a real
trade-off patient from one who landed there by chance. The correction is a
nonparametric bootstrap done at the right level: resample patients with
replacement, \emph{refit} the tuned forest on each resample, and predict on the
original unshuffled covariates, so that each patient keeps their own row.
Fifty replicates give a 95\% interval for every patient's position on both
axes. The result is in Figure~\ref{fig:tradeoff-bootstrap}, and it is
categorical: \textbf{0 of 4,579} patients have their whole confidence region
inside the bottom-right quadrant, against 1,544 whose point estimate is there.
On both axes the within-patient bootstrap standard deviation is roughly twice
the across-patient spread of the estimates — a signal-to-noise ratio of $0.57$
on the bleeding axis and $0.52$ on the ischemic one — which is the same
statement as the placebo floor of the previous section, arrived at by
resampling rather than by permutation.

\begin{figure}[h]
\centering
\includegraphics[width=1\textwidth]{CATE_tradeoff_bootstrap.png}
\caption{The trade-off plane with bootstrap uncertainty, from 50 refits of the
tuned forest on resampled cohorts. Left, per-patient 95\% intervals for a
random subsample of 1,000 patients; right, the whole cohort, with patients
whose interval lies entirely in the trade-off quadrant marked in red. There
are none.}
\label{fig:tradeoff-bootstrap}
\end{figure}

What \emph{does} survive the bootstrap is the average: 4,541 of 4,579 patients
(99.2\%) have a significantly positive bleeding benefit of abbreviating. The
plane is thus a precise picture of the whole chapter — the effect is real,
strong and essentially the same for everyone.

\paragraph{A caution about the conflict subgroup.}
Flagging the top 15\% by conflict score produces a subgroup of 687 patients
whose \emph{observed} event rates look spectacular: 24.6\% bleeding under
prolonged therapy against 3.7\% under abbreviated, a raw difference of 20.8
percentage points, against 1.6 points in the rest of the cohort. The number is
real and it means nothing, because the score that selected those patients was
computed from a model fitted on the same patients and the same outcomes. The
subgroup is selected \emph{to} exhibit that contrast; quoting its observed
event rates as an effect estimate would be circular. It is reported here as a
descriptive, in-sample profile only, and the honest counterpart is the previous
paragraph: with error bars, none of these patients is distinguishable from the
cohort average.

\section{From Ranking to Decision}
\label{sec:decision}

Everything up to here has tested a statement about a model — that its ranking
of patients tracks real effect differences. A weaker and more defensible
question remains, and it is the one a clinician would actually ask: not whether
the ranking is good, but whether \emph{acting} on it produces a measurable
gain. A rule can be valuable even with a mediocre ranking, and a good ranking
can be worthless if the gain it concentrates is smaller than the cost of
deviating from the simple strategy. The two questions are answered with
different machinery, so the negative result of the previous sections does not
settle this one by itself.

\paragraph{The net-benefit endpoint.}
A decision needs an outcome that already encodes the trade-off, since no single
endpoint can. This section therefore adds the trial's own composite: NACE (Net
Adverse Clinical Events, 7.7\% in this cohort), which combines death, MI,
stroke and BARC 3--5 bleeding, alongside the ischemic composite MACCE (6.0\%).
All estimates are cross-fitted AIPW with the known randomization propensity.

The sanity anchor holds first: the doubly-robust ATE on NACE is $+0.005$ with
a 95\% interval of $[-0.011,\,+0.020]$ — net-neutral, spanning zero, which is
the trial's primary result recovered by the same pipeline that will be asked to
evaluate the rules. On the components the same pipeline gives $+0.027$
$[+0.011,\,+0.042]$ for BARC 2/3/5 bleeding and $-0.000$ $[-0.014,\,+0.013]$ for
MACCE: the bleeding benefit of abbreviating is there, the ischemic cost is not
measurable, and their composite is a wash.

\paragraph{Pre-specified subgroups.}
Before any learned rule, the higher-powered test: seven pre-specified effect
modifiers — high PRECISE-DAPT score, acute coronary syndrome presentation,
diabetes, eGFR $< 60$, age $\geq 75$, anemia, oral anticoagulation — each
evaluated as the difference between the doubly-robust ATE inside and outside
the subgroup. A low-dimensional pre-specified test of this kind is far more
powerful than the omnibus search already performed, which is why it is worth
running even after a null. Across the three endpoints reported (NACE, BARC
2/3/5 bleeding, MACCE) and the seven modifiers, \emph{every one} of the 21
interaction intervals contains zero. The largest point difference is on anemia
for NACE, $+0.026$ with an interval of $[-0.041,\,+0.092]$, which is to say the
data cannot distinguish it from nothing on either side.

\paragraph{Policy value.}
The crux. Two learned rules — an interpretable doubly-robust policy tree and a
policy forest — are compared against the two fixed strategies on NACE, with
honest cross-fitted values and a paired interval on the improvement
(Table~\ref{tab:policy-value}).

\begin{table}[h]
\centering
\caption{Cross-fitted AIPW policy value on NACE (lower is better) and paired
improvement over the better fixed arm. A usable personalized rule would show a
positive improvement with an interval excluding zero.}
\label{tab:policy-value}
\small
\begin{tabular}{lccc}
\toprule
Policy & Value (NACE) & \% prolonged & Improvement vs.\ best fixed [95\% CI] \\
\midrule
All abbreviated   & 0.076 & 0.0   & --- \\
All prolonged     & 0.080 & 100.0 & $-0.005$ $[-0.020,\,+0.011]$ \\
ATE sign          & 0.087 & 40.0  & $-0.011$ $[-0.021,\,-0.002]$ \\
Policy tree       & 0.082 & 31.1  & $-0.006$ $[-0.015,\,+0.002]$ \\
Policy forest     & 0.087 & 31.6  & $-0.011$ $[-0.019,\,-0.003]$ \\
\bottomrule
\end{tabular}
\end{table}

No rule beats \emph{abbreviate for everyone}. Two of the three learned rules
are significantly \emph{worse} than it, which is the expected cost of
personalizing on a signal that is not there: every patient the rule diverts to
prolonged therapy is a patient who, on average, gains nothing and bleeds more.

\paragraph{The trade-off frontier is empty at every weight.}
One objection survives the table above: NACE weights its components equally,
and a clinician who weighted ischemic events more heavily might find a rule
worth having. The objection is testable. Sweeping a weighted outcome
$Y_w = w \cdot \text{MACCE} + (1-w) \cdot \text{BARC 3--5}$ from $w = 0$ to
$w = 1$ moves the better fixed arm from all-abbreviated to all-prolonged, as it
should — the trade-off is real at the population level. At no weight does a
learned rule beat the better fixed arm: the improvement is $-0.001$ at $w = 0$,
$-0.005$ at $w = 0.5$ and $-0.003$ at $w = 1$, negative throughout, and at
$w = 0.5$ the interval excludes zero on the wrong side. There is no clinical
trade-off preference, however extreme, at which personalization on this cohort
pays.

\section{How Small a Heterogeneity Would Have Been Seen}
\label{sec:delta}

The positive control of Section~\ref{sec:positive-control} establishes that the estimators recover an
average effect. It does not establish that they would have recovered
heterogeneity of any particular magnitude, and these are different claims: the
second is the one a negative result actually needs. Without it, «no
heterogeneity was found» is an assertion about the analysis and not about the
data. This section replaces the assertion with a number, one per endpoint: the
smallest heterogeneity $\delta$ that the protocol used in this chapter would
reliably have detected.

\paragraph{Protocol.}
The cohort's real covariates and real treatment assignment are left untouched;
only the outcome is resampled. A baseline risk $\mu(x)$ is fitted on the
abbreviated arm, a direction of heterogeneity $g(x)$ is fixed as the
standardized patient age — chosen because it is the second most important
covariate in the already-fitted bleeding model, and because age is a recognised
component of standard bleeding-risk scores, so the injected modifier is
clinically plausible rather than arbitrary — and synthetic outcomes are drawn
as $Y \sim \mathrm{Bernoulli}\big(\mathrm{clip}(\mu(x) + T \cdot \delta \cdot
g(x))\big)$. For the two bleeding endpoints $\delta$ is an absolute risk
difference; for the three rare endpoints an additive $\delta$ of comparable
size would distort the prevalence, so it is applied as a log odds ratio on the
logit scale. The tuned causal forest is then refitted with unchanged
hyperparameters and scored with the same cross-validated AUTOC lower bound used
everywhere else.

The comparison at each $\delta$ is not against zero but against a noise floor
recomputed \emph{at that same} $\delta$ by permuting the treatment: the same
synthetic outcomes, the same estimator, no effect. This matters because the
floor is not constant — a rarer or more skewed outcome shifts it — and
comparing against a fixed zero would credit the protocol with detections it did
not make. With 20 seeds per $(\text{endpoint}, \delta)$ pair and 500 fits in
total, $\delta^*$ is defined as the smallest $\delta$ on the grid at which the
observed score beats its paired floor in at least 90\% of seeds.

\begin{table}[h]
\centering
\caption{Minimum detectable heterogeneity $\delta^*$ per endpoint, against the
endpoint's own average effect. Detection requires beating the permuted-treatment
floor in $\geq 90\%$ of 20 seeds.}
\label{tab:delta-star}
\begin{tabular}{lrrl}
\toprule
Endpoint & Real ATE & $\delta^*$ & Reading \\
\midrule
Bleeding             & $+0.045$ & $0.05$ risk diff. & as large as the entire average effect \\
Bleeding BARC 2/3/5  & $+0.028$ & $0.05$ risk diff. & larger than the average effect \\
MI                   & $-0.005$ & $\mathrm{OR} = 3.0$ & clinically implausible \\
Cardiovascular death & $+0.003$ & $\mathrm{OR} = 5.0$ & clinically implausible \\
Stroke               & $+0.005$ & not reached & degenerate, see below \\
\bottomrule
\end{tabular}
\end{table}

\begin{figure}[h]
\centering
\includegraphics[width=1\textwidth]{delta_sensitivity.png}
\caption{Observed AUTOC lower bound (blue) against the permuted-treatment floor
(grey) as the injected heterogeneity grows, one panel per endpoint; bands are
the 10th--90th percentile across 20 seeds and the dashed line marks
$\delta^*$. The two curves separate cleanly once the injected signal is large
enough, which is what makes the protocol a measurement rather than a
formality.}
\label{fig:delta-sensitivity}
\end{figure}

\paragraph{Reading the number honestly.}
The protocol works: Figure~\ref{fig:delta-sensitivity} shows the observed curve
lifting away from the floor as soon as the injected heterogeneity is large
enough, on four of five endpoints. That is the property the positive control
could not supply — evidence that this pipeline detects heterogeneity when
heterogeneity is present.

But $\delta^* = 0.05$ on bleeding \emph{is large}. It is the entire average
effect of the trial, expressed as variation around it. What these data exclude
is heterogeneity of that magnitude, not heterogeneity as such: at
$\delta = 0.02$ — half the average effect, and clinically relevant by any
standard — the protocol detects in 80\% of seeds, below the 90\% threshold, and
at $\delta = 0.01$ in 50\%. This is precisely why the conclusion of this thesis
is «it was not found» and never «it does not exist», and it is worth stating
before a reader concludes it for themselves.

On the ischemic endpoints the reading is different in kind rather than in
degree. An odds ratio of 3 to 5 for effect modification is not a magnitude
anyone expects in cardiology; asking for it is asking for something the cohort
was never going to supply. The correct claim on MI and cardiovascular death is
therefore not «there is no heterogeneity» but \textbf{«it is not measurable on
this cohort»}, and the $\delta^*$ column turns that from a hedge into a
measurement. Chapter 7 takes up the reason.

\paragraph{Two defects of this analysis, stated plainly.}
First, the baseline $\mu(x)$ underestimates prevalence: it generates 8.5\%
synthetic bleeding against 11.1\% real, and 6.2\% against 7.8\% on BARC 2/3/5.
Fewer synthetic events mean less power, so the $\delta^*$ values in
Table~\ref{tab:delta-star} are \emph{conservative} — a correctly calibrated
baseline would move them down, not up. This does not weaken the negative
result, but it does mean the numbers are upper bounds on the true minimum
detectable effect.

Second, the stroke row is not interpretable and is reported as such. Under the
null, where the observed score should beat its floor in about half the seeds by
construction, it does so in 85\% — the protocol is broken on that endpoint
before any heterogeneity is injected. The cause is visible in the same output:
$\mu(x)$ generates roughly 14 synthetic strokes against 35 real ones, and at
that event count the AUTOC and its floor are both essentially degenerate.

\todo[inline]{Before the final draft: recalibrate $\mu(x)$ so synthetic
prevalence matches the real event rates, redo the \texttt{stroke} row, and
extend the protocol from the causal forest to the other four estimators.}

\section{What This Chapter Establishes}

Four things, in the order they were established. The average effect of MASTER
DAPT is recovered by five independent estimator families and matches the
published trial, so the tools work. The variation of that effect across
patients does not exceed the level that the same tools produce from permuted
data, does not survive a doubly-robust held-out evaluation under either
weighting, and grows with model capacity in the way a fitting artefact does and
a data property does not. The trade-off plane that a personalized rule would
need is a single compact cloud with no patient whose position in the decisive
quadrant survives a bootstrap. And the decision-level question, asked directly
with cross-fitted policy values, returns the same answer: nothing beats
abbreviating for everyone, at any weighting of the bleeding--ischemia trade-off.

The sensitivity analysis bounds all of this. On the bleeding endpoints, the
protocol would have caught heterogeneity as large as the trial's own average
effect, and there is none; on the ischemic endpoints, it would have needed
effect modification of a magnitude nobody claims to exist, so the honest verdict
there is that the question is not answerable on this cohort at all.

What the chapter does not do is explain \emph{why}. A negative result without a
mechanism is one bad-luck cohort away from being uninformative, and the two
observations most in need of an explanation are already on the table: the
cohort's bleeding risk was flattened before any model saw it — the same fact
that Chapter 5 left open when bleeding, with 506 events, was predicted worse
than cardiovascular death with 81 — and the ischemic composite sums endpoints
whose effects point in opposite directions. Chapter 7 takes both.
